\section{Environment}
\label{sec:environment}

We release \textsc{PokeGym}, a Gymnasium-compatible
environment~\citep{gymnasium} that wraps a headless
PyBoy~\citep{pyboy} emulator running an unmodified Pok\'emon Red ROM.
The environment exposes a discrete action space, three first-class
observation representations, an event-flag-based reward calculator, and
a curriculum-aware save-state system.

\subsection{Action Space}
The base environment supports the standard Game Boy buttons: D-Pad
(\textsc{Up}, \textsc{Down}, \textsc{Left}, \textsc{Right}),
\textsc{A}, \textsc{B}, \textsc{Start}.  We exclude \textsc{Select}
because it has no in-game effect for the section of the game we
study.  This yields $|\mathcal{A}| = 7$.  Each environment step
advances the emulator by 24 frames (one in-game tile of overworld
movement) before returning control to the agent.

\subsection{Observation Treatments}
\label{sec:obs-treatments}

The same emulator state can be exposed as three different observation
spaces, selected at environment construction time via
\texttt{observation\_type}.

\paragraph{Pixel.}
A single $80 \times 72$ grayscale frame ($\mathcal{B}^{80 \times 72}$,
$\mathcal{B} = \{0, \ldots, 255\}$).  We use grayscale rather than RGB
because the Game Boy's 4-shade palette already loses no information
relative to RGB, and this halves the input dimensionality.  The pixel
agent encodes this with a Nature-DQN-style
CNN~\citep{mnih2015atari}---three convolutional layers followed by a
fully-connected projection to a 256-dim feature vector.

\paragraph{Symbolic.}
A flat 29-dimensional float vector containing:
%
\begin{itemize}
\item \textbf{Player position} (3 dims): map ID, $x$ tile coordinate,
  $y$ tile coordinate.
\item \textbf{Health and party} (4 dims): current HP fraction, party
  size, max party level, total levels summed.
\item \textbf{Currency and inventory} (3 dims): money (log-scaled),
  badge count, item-bag fill ratio.
\item \textbf{Exploration state} (2 dims): unique-maps-visited count,
  unique-tiles-visited count.
\item \textbf{Event flags} (17 dims): a binary vector for the 17
  pre-registered Boulder Path event flags (\Cref{sec:events}); the
  18th flag (\texttt{EVENT\_GOT\_POKEMON\_FROM\_FAN\_CLUB\_CHAIRMAN})
  is held out as a sanity-check canary.
\end{itemize}
%
The symbolic agent encodes this with a 2-layer MLP ($29 \to 128 \to
256$) projecting to the same 256-dim feature space as the pixel CNN,
keeping the downstream LSTM input dimension matched.

\paragraph{Hybrid.}
A \texttt{Dict} space combining both branches.  The CNN and MLP run in
parallel; their 256-dim outputs are concatenated to a 512-dim vector
before entering the LSTM.  The hybrid agent has roughly 2$\times$ the
encoder parameters of pixel or symbolic alone, but identical
recurrent and policy heads.  We do not match parameter counts across
treatments---matching the encoder \emph{capacity} would require
shrinking the pixel CNN below the published Nature-DQN architecture
and would conflate two design decisions; we match the architecture
above the encoder instead.  We discuss this trade-off in
\Cref{sec:discussion}.

\subsection{Reward}
\label{sec:reward}

We use an event-flag-based reward calculator that issues a fixed
positive reward the first time each pre-registered flag is triggered
in an episode, plus a small per-step time penalty.  The 18-flag set is
committed in the pre-registered plan and corresponds to the canonical
critical path between the player's bedroom in Pallet Town and Brock's
Boulder Badge in Pewter City Gym.  In addition to event flags, we
issue:
%
\begin{itemize}
\item A small per-step time penalty ($-0.001$) to encourage timely
  progress.
\item A discovery reward for the first visit to each new map ID.
\item Negative reward for fainting (HP $\to$ 0 across the entire
  party).
\end{itemize}
%
All other shaping signals (e.g., level-based or money-based shaping)
are zero by default.  This keeps the reward signal close to the
event-flag milestones, which are observable, interpretable, and
directly map to the game's progression.

\subsection{Pre-Registered Event Flags}
\label{sec:events}

The 18 event flags used by the reward calculator are listed in the
appendix and committed in the pre-registered plan.  They cover:
deliver Oak's parcel (8 flags), traverse Viridian Forest (3 flags),
defeat Brock's gym trainers and Brock himself (5 flags), plus one
optional stretch milestone (Rival fight on Route 22) and one canary
flag that should not be reachable before Brock and is monitored as a
sanity check.

Flag bits and addresses are extracted from the
\texttt{pret/pokered}\footnote{\url{https://github.com/pret/pokered}}
disassembly and verified by emulator-side replay.

\subsection{Curriculum and Episode Resets}
\label{sec:curriculum}

We use a single fixed save state, \texttt{post\_intro.state}, captured
in Pallet Town immediately after Professor Oak's introduction
sequence.  All episodes reset to this state.  Pilot results in this
paper therefore reflect the agent's ability to traverse Route 1 →
Viridian → Viridian Forest → Pewter from a fixed starting point.
Future work will sample from a curriculum of save states (Pallet,
Viridian Pok\'emon Center, Mt. Moon entrance) using a Go-Explore-style
weighted distribution.

\subsection{Episode Truncation}
Episodes truncate at 15{,}000 environment steps regardless of in-game
state.  This is approximately 30 minutes of in-game time at the
emulator's 24-frame action repeat and is empirically sufficient for
even a randomly-acting policy to leave Pallet Town occasionally,
preventing the agent from being stuck in a state with zero reward
signal.

\subsection{Reproducibility}
\textsc{PokeGym} is open-source and pip-installable.  The training
script, evaluation harness, statistical analysis pipeline, and live
training dashboards are released under an OSI-approved license at the
project repository.  We seed every RNG that affects training (Python,
NumPy, PyTorch, Gymnasium) deterministically per run and document the
exact seeds in the supplementary material.
